\begin{table}[H]
\centering
\caption{Marginal Return to Assignment Value by Delivery Fidelity}
\label{tab:struct_signal_delivery_margins}
\begin{threeparttable}
\scriptsize
\setlength{\tabcolsep}{3pt}
\begin{tabular}[t]{>{\raggedright\arraybackslash}p{3.3cm}ccccc}
\toprule
Scenario & $\omega$ & $\tau$ & Gain $G_N\to G_K$ & MP of $G$ & Cross-partial \\
\midrule
Nigeria realized & 0.84 & 0.31 & $+0.001$ & 0.002 & 0.023 \\
Nigeria high delivery & 0.84 & 0.90 & $+0.113$ & 0.207 & 1.061 \\
Nigeria high delivery + execution & 0.95 & 0.90 & $+0.128$ & 0.234 & 1.198 \\
Kenya observed delivery & 1.00 & 0.50 & $+0.009$ & 0.016 & 0.151 \\
Kenya high delivery & 1.00 & 0.90 & $+0.134$ & 0.246 & 1.261 \\
Fully high-input & 0.98 & 0.95 & $+0.169$ & 0.309 & 1.503 \\
\bottomrule
\end{tabular}
\begin{tablenotes}[para,flushleft]
\footnotesize
\item \textit{Notes:} The table evaluates the assignment-payoff term $\hat{\lambda}G\omega\tau^{\hat{\alpha}}$. ``Gain'' is the ATE increase from raising assignment value from Nigeria's estimate to Kenya's estimate, holding the row's $\omega$ and $\tau$ fixed. The last two columns report the analytic marginal product of assignment value and the assignment-value--delivery cross-partial implied by the preferred production map.
\end{tablenotes}
\end{threeparttable}
\end{table}
